\FloatBarrier
\section{Open problems}
\label{sec:open}

The experiments answered some questions and opened others. Each problem below comes from an empty or weak cell of Table~\ref{tab:dictionary}, from a result of Section~\ref{sec:experiments}, or from a place where two clocks meet, and each is stated with a first test. Problems \op{9} and \op{10} are conjectures that come from the mechanics of nonlinear oscillators; we include them because they make predictions that can be checked, and we do not claim evidence for them yet.

\begin{enumerate}[label=\textbf{OP\,\arabic*.}, leftmargin=*, itemsep=3pt]
\item \textbf{Coupling of clocks 1 and 2.} Which features of the Lyapunov profile of the forward pass settle before the loss does? In two small models the whole-stack exponent settles within 3\% of training while the per-layer exponents keep growing (Experiment~1). The next test is the full Jacobian spectrum of every layer across the public checkpoints of a model family above a billion parameters.

\item \textbf{Periodic orbits in depth.} Do models that reuse a block many times, such as looped and recurrent-depth models \citep{prairie2026parcae}, have periodic or quasi-periodic orbits in the residual stream, and do these orbits carry computation? Return maps of hidden states over many loop iterations would show them directly.

\item \textbf{The full Lyapunov spectrum of a language model.} How many directions expand along a forward pass and along a generation, and what is the resulting attractor dimension? QR-based algorithms \citep{engelken2020lyapunov,vogt2020lyapunov} make this feasible for models of GPT-2 size, and the exponent of clock~3 would turn exposure bias into a measured quantity.

\item \textbf{The type of a loss spike.} Are the loss spikes of pretraining crossings of the stability boundary, period-doubling cascades, catapults \citep{lewkowycz2020large} or escapes from a basin? Each type leaves a different signature in the curvature and in the update direction just before the spike, so open training logs and checkpoints can be classified by these signatures.

\item \textbf{Escape laws across samplers and scales.} How do nucleus sampling and repetition penalties change the barrier of Proposition~\ref{prop:arrhenius}, and do the saturated barrier and the locking temperature keep growing with model size beyond 1B parameters? Reasoning models add the recurrent orbits of Experiment~3, which have no period: do they have an escape law of their own, and do recurrence plots of the trace predict when a model will run out of budget? A harder version asks whether the onset of a loop and the coherent window of a prompt can be read from the prompt and the weights before any text is generated.

\item \textbf{Fractal trainability of transformers.} Does the fractal boundary between convergent and divergent hyperparameters \citep{sohldickstein2024boundary}, reproduced for a small network in Figure~\ref{fig:basins}, persist for transformers trained with Adam, and how does its dimension depend on depth? Fine grids over learning rate and warmup length give a first estimate.

\item \textbf{Metastability in generation.} Is the length of a coherent stretch of generated text before drift a metastable residence time, exponentially long in some barrier, like the cluster states in depth \citep{geshkovski2024dynamic}? Survival curves of coherence against temperature would test this.

\item \textbf{Koopman modes of a forward pass.} Can a linear operator fitted to residual-stream trajectories predict the effect of an intervention several layers ahead, and do its eigenfunctions align with known features? Dynamic mode decomposition on residual streams, compared with the tuned lens \citep{belrose2023eliciting}, is a direct test.

\item \textbf{Targeted energy transfer (conjecture).} In coupled mechanical oscillators, a nonlinear energy sink absorbs vibration energy from a primary system in one direction and irreversibly \citep{vakakis2009nonlinear,farid2016tuned}. We conjecture a similar one-way transfer of signal energy between spectral modes of the residual stream, with massive activations and attention sinks \citep{sun2024massive,xiao2023efficient} in the role of the sink. Tracking the energy of each singular mode of the layer Jacobians across depth and training would confirm or refute it.

\item \textbf{Escape under internal forcing (conjecture).} At the edge of stability the iterate oscillates along stiff directions, and the other directions see this oscillation as periodic forcing. A periodically forced particle escapes from a potential well when the forcing amplitude crosses a threshold that depends on the forcing frequency \citep{farid2019escape,farid2022escape}, and we conjecture that some loss spikes are escapes of this kind. Gradient descent on a quartic well coupled to an oscillating direction is the simplest test. Experiment~2 adds a related question: self-stabilization \citep{damian2022selfstabilization} explains a single mode at the edge, but our transformer has several, which exchange order at every step, and their averaged dynamics is open.

\item \textbf{Early warning signals.} Near a bifurcation, recovery from perturbations slows down and fluctuations grow. Do grokking, induction-head formation and loss spikes show rising autocorrelation and variance of the loss or of the gradient norm before they happen? Runs with a known transition step \citep{power2022grokking,olsson2022incontext} allow a direct check.

\item \textbf{Control of clock 3.} Can a controller built from the escape analysis keep generation coherent at a lower cost than repetition penalties, with Mirostat \citep{basu2020mirostat} and truncation sampling \citep{hewitt2022truncation} as baselines? Three first tests on GPT-2 (Appendix~\ref{app:exp3}) fail in an instructive way. Heating the sampler when a loop is detected cuts looping generations at $T=0.3$ from 68\% to 6\%, but 94\% of the outputs remain repetitive. Rolling the text back to the first token of a loop and resampling there removes all exact loops at $T=0.5$, but 91\% of the outputs stay repetitive. A feedback controller on recurrence and surprisal settles near the center of the coherent window from any start, but its share of coherent texts (21\% to 83\%) does not beat the 82\% of the best fixed temperature. All three act on the noise, while the trap is stored in the context. Controllers that act on the state, by editing the cache or by adding a steering vector to the residual stream \citep{turner2023steering,zou2023representation}, are the natural next step: in mechanical terms, a control force on the state instead of a change of the noise intensity.
\end{enumerate}
